\begingroup
\renewcommand{\arraystretch}{1.45}
\fancyhead[C]{}
\clearpage\fancyhead[L]{\small REMARKS.}\pdfbookmark[0]{Remarks}{final-bemerkungen}\SourceLeaf{500}
\begin{center}\Large\textbf{REMARKS.}\end{center}
\begin{center}\rule{.09\textwidth}{.4pt}\end{center}

In this third volume of \textsc{Gauss}’s works I have brought together all the treatises and papers from general analysis, specifically from the theory of algebraic functions, Gaussian series and elliptic functions, together with some communications concerning \textsc{Pfaff}’s theorem and the construction of logarithm tables. They comprise a doctoral dissertation (previously printed in quarto); previously published treatises: five in the \textit{‘Commentationes societatis regiae scientiarum Gottingensis’} (quarto), one in the \textit{‘Abhandlungen der Königlichen Gesellschaft der}
\textit{Wissenschaften zu Göttingen’} (quarto) and one in \textsc{Crelle}’s \textit{‘Journal für reine und angewandte Mathematik’} (quarto); also six notices of his own treatises in the \textit{‘Göttingischen Gelehrten Anzeigen}
\textit{und Nachrichten’} (octavo); twelve communications on writings by others in the \textit{‘Göttingischen gelehrten Anzeigen’} (octavo) (unsigned by \textsc{Gauss}, but his authorship confirmed by the records of the Göttingen University Library), in the \textit{‘Monatlichen Correspondenz zur Beförderung der Erd-}
\textit{und Himmels-Kunde herausgegeben vom Freiherrn \textsc{von Zach}’} (octavo), and in the \textit{‘Astronomischen Nachrichten}
\textit{herausgegeben von \textsc{Schumacher}’} (quarto); and also two explanations, in the \textit{‘Monatlichen Correspondenz’} and in \textsc{Vega}’s \textit{‘Sammlung von Hülfstafeln’} (quarto), of the Gaussian addition and subtraction logarithm tables. I have not had these tables themselves reprinted here, since they are very widely available and the format of these works would be inconvenient for using such tables.

In editing I have followed the same principles as in the earlier volumes. To make the subjects in so extensive a volume easier to survey, the principal theorems are distinguished typographically in the same way. To make both the earlier editions and the present one easier to use, in references to writings not contained in this volume I have given their own titles instead of the place of publication, and the article numbers instead of the page numbers. On page 20, in lines 15 and 16, in accordance with a handwritten remark, \textit{‘objectionem secundam et quar\SourceLeaf{501}%
tam’} has been put instead of \textit{‘objectionem tertiam et quartam’}, and in line 19 likewise \textit{‘objectionem tertiam’} instead of \textit{‘objectionem}
\textit{primam’}. Otherwise the present edition differs from the earlier editions of the same writings only by the correction of some printing errors, for which I was able in part to use \textsc{Gauss}’s original manuscripts.

From the manuscript estate I have included the notes attached to pages 30 and 112, the second part of the treatise \textit{‘Disquisitiones generales circa seriem infinitam etc.’}, a detailed treatise on interpolation with some additions explained in my remarks on page 328 and following, and a series of treatises concerning elliptic functions.

Here, as previously, I have had the manuscript records printed without alteration apart from the correction of insignificant slips of the pen, and have marked off my insertions by [brackets]. With the exception of the \textit{‘Fortsetzung der Untersuchungen über das arithmetisch-geometrische Mittel’}, these insertions consist only of short sentences. I have brought together the historical information and necessary additions concerning the posthumous papers on hypergeometric series and interpolation in the remarks immediately following those treatises.

\bigskip
The Royal Society of Sciences at Göttingen wished \textsc{Riemann} to edit \textsc{Gauss}’s works on elliptic functions, and the manuscripts were handed over to him for that purpose. Unfortunately he left neither a written nor an oral communication from these studies. Only after \textsc{Riemann}’s regrettably early death, and after this third volume had been printed apart from that portion, was I able to take over the Gaussian manuscripts for my own editorial work.

\bigskip
\textsc{Gauss} published only part of his investigations of the functions we now call elliptic: an application of this theory to higher arithmetic in the \textit{Summatio quarundam serierum singularium}, 1808 \textit{September}, and an application to the determination of the secular perturbations of the planets in the \textit{Determinatio attractionis, quam in punctum quodvis positionis datae exerceret planeta, si}
\textit{ejus massa per totam orbitam ratione temporis, quo singulae partes describuntur, uniformiter esset dipertita,} 1818 \textit{January.} He set down the greater part in some incomplete drafts of the beginnings of various treatises and in numerous records of individual formulas scattered very widely among other work. I have arranged these investigations from the manuscript estate here in separate groups, according as they start from the algorithm of the arithmetic–geometric mean, or from another algorithm analogous to it and connected with it, or concern the special lemniscatic functions, or, finally, use representations of the general functions by products as an essential aid. The investigation of the \textit{Pentagramma}
\textit{mirificum} is an application of the division of complete elliptic integrals into five parts.

\bigskip
The treatise with the two sections \textit{de origine proprietatibusque generalibus numerorum mediorum arithmetico-geometricorum} and \textit{de functionibus transcendentibus quae ex differentiatione mediorum}
\providecommand{\TStart}{\hspace*{.5em}}\providecommand{\TCont}{\hspace*{1em}}\providecommand{\TD}{\hspace{.3em}.\hspace{.3em}.}
\SourceLeaf{502}
\textit{arithmetico-geometricorum oriuntur} forms the first entry, and a very carefully written one, in a handbook which, as its title states, was used from the year 1800 onwards. Work on this subject will, however, have begun at a much earlier time; according to reports of an oral statement by \textsc{Gauss}, he seems to have known in 1794 the relations between arithmetic–geometric means and the power series whose exponents advance through the square numbers.

The formulas for the algorithm included in Art. 18, page 389, which reduces the new transcendental functions introduced by \textsc{Gauss}, with the squared values of both arguments, to the transcendental functions with the simple argument values, follow immediately in a handbook after an astronomical calculation ending with the words `\textit{completed on 2 May} 1809'. The records of the other investigations of the arithmetic–geometric mean occur partly on individual undated sheets and partly as short entries in the handbooks at places earlier left empty between different kinds of work to make them easier to survey; they permit no certain dating.

As already remarked in Art. 12, I thought it necessary, for the reader’s convenience, to unite these very fragmented investigations in a connected exposition, even at the risk of setting out here some developments which \textsc{Gauss} did not carry out. These include, for example, the derivation of the differential equation for the arithmetic–geometric mean without presupposing the expansion in series and the representation by definite integrals, a derivation which follows on from the investigation carried out in Art. 10 and differs from the first derivation of this kind given by Mr. \textsc{Borchardt}. The question raised in Art. 17 concerning the connection between binary quadratic forms with negative determinants and the new transcendental functions found by \textsc{Gauss} receives its most complete resolution in Mr. \textsc{Kronecker}’s investigations of this subject.

\medskip
For the \textit{lemniscatic functions} we possess the date recorded by \textsc{Gauss} in a handbook: `\textit{We had begun to consider lemniscatic functions} 1797. \textit{January.} 8.' Judging by the paper and the form of the handwriting, the articles of the \textit{first group} of investigations of lemniscatic functions, written on a separate sheet and designated here by me as 1, 2, 3, 4, seem to be the earliest of the records in the estate. The following records of the same first group, assembled here on pages 406--412 under the common heading used by \textsc{Gauss} in several places, `\textit{Einige neue Formeln die Lemniscatischen Functionen betreffend}', belong to a much later time. Their results were probably found much earlier and written down anew, partly as a test of memory and partly in the pursuit of especially elegant formulas, in two handbooks, one of which was begun in `\textit{November} 1801', the other `\textit{in May} 1809 \textit{begun}'. \textsc{Gauss} everywhere denotes the functions sin. lemn. and cos. lemn. by characters formed by contracting \(s\) and \(l\), and \(c\) and \(l\), respectively. Since these characters are not yet available in printing type, I have replaced them here by the words themselves.

\SourceLeaf{503}

Of the investigations chiefly concerning the representation of lemniscatic functions by infinite products and trigonometric series, which I have assembled as a \textit{second group}, the first four articles form the \textit{Scheda prima} of a notebook begun in July 1798 according to its title page. The articles 1, 2, 3 were numbered by \textsc{Gauss} himself; then, in the same notebook, the contents of Art. [4.] [5.], page 415, immediately follow, and after these, with frequent interruptions by astronomical investigations, the portion of the calculations included in Art. [17.], page 431, which is explained in Latin, and the separate portions of the contents of Art. [6.] [7.], page 417,418. The contents of Art. [8.] [9.] [10.], pages 418, 419, and Art. [15.] [16.], pages 423--425, are scattered through a notebook begun in 1799 November. The records of the remaining investigations probably belong to a later time, perhaps with the exception of the division of the whole lemniscatic arc into five parts in Art. [13.]. Apart from the summation formulas for the lemniscatic integral of the second kind reproduced in Art. [14.] and entered in a handbook after 20 February 1817, they are all on separate sheets.

\medskip
The earliest record concerning general elliptic functions preserved in the manuscript estate is probably the treatise \textit{Zur Theorie der neuen Transscendenten} I, begun in November 1799 and printed on pages 433--435. It is scattered among investigations of quite different subjects in a notebook whose title page bears the inscription `\textit{Varia, imprimis de Integrali} \(\displaystyle\int\frac{\mathrm{d}u}{\surd(1+\mu\mu\sin u^{2})}\), \textit{Novembr.} 1799'.

\medskip
Treatise II, headed `\textit{Zur Theorie der transscendenten Functionen gehörig}', follows immediately in a handbook after the auxiliary tables reproduced in the \textit{Theoria motus corporum coelestium}. These investigations, particularly those concerning division into seven parts, probably belong to the year 1808 and will have occasioned a communication to \textsc{Schumacher}.

\medskip
Treatise III, with the opening words `We arrange the theorems concerning those series and infinite products which belong to the theory of arithmetic–geometric means as follows:', pages 446--460, follows immediately in a handbook after an astronomical calculation to which the remark `completed on 28 April 1809' is attached.

\medskip
Treatise IV, `Hundert Theoreme über die neue Transscendente', pages 461--469, is on individual sheets without any indication of date.

\medskip
Treatise V, pages 470--480, with the opening words `General Theorem', contains the two dates 1827 Aug. 6 and Aug. 29. The last part of this paper probably forms one of \textsc{Gauss}’s earliest preliminary studies for his `Allgemeinen Lehrsätze in Beziehung auf die im verkehrten Verhältnisse des Quadrats der Entfernung wirkenden Anziehungs- und Abstossungskräfte', and suggests that he had recognized the connection of this field with another field of analysis, namely the \SourceLeaf{504}theory of functions with complex arguments, a connection which led \textsc{Riemann} to so fruitful a field of research.

\medskip
The investigation of the \textit{Pentagramma mirificum}, a spherical pentagon whose five diagonals are quadrants, occurs at two separate places in a handbook, with Art.\,[4.] on a separate sheet. What I have designated here as Art.\,[1.] and [2.] was written before 23 January 1836; the other part contains in Art.\,[7.] the date 1843 April 20.

\medskip
Among the investigations of the new transcendental functions assembled under No.\ I there is also the beginning of a treatise headed ‘\textit{Motus solidi a nullis viribus sollicitati}’. The problem is carried there as far as the elliptic integral appearing in the known solution, so that it may be supposed that \textsc{Gauss} had recognized that the equations determining this motion can be expressed in finite form with the aid of the new transcendental functions.

The record of the investigations of the new transcendental functions assembled here under V, with the dates 1827 August 6 and August 29, was probably occasioned by \textsc{Jacobi}’s letters to \textsc{Schumacher}, dated at Königsberg 1827 June 13 and Aug.\ 2. The first gives the algebraic equations for division of elliptic integrals into three and five parts, and the other the equation between the trigonometric tangents of the arguments for transformation of arbitrary degree. These two letters were published in September 1827 in No.\ 123 of the ‘Astronomischen Nachrichten’, but the originals had previously been communicated to \textsc{Gauss}, as is evident from \textsc{Gauss}’s letters to \textsc{Schumacher} of 4 and 19 Aug.\ 1827. \textsc{Jacobi} himself published the proofs of those theorems of his in a letter to \textsc{Schumacher} dated at Königsberg 13 November 1827 and printed in December 1827 in No.\ 127 of the ‘Astronomischen Nachrichten’.

It was not yet at that time, but later in the same year, that \textsc{Abel}’s treatise ‘\textit{Recherches sur les fonctions elliptiques} \S.\ I.\,---\,\S.\ VII.’ appeared in the second issue of the second volume of \textsc{Crelle}’s Journal für reine und angewandte Mathematik. This was probably the work of \textsc{Abel} about which \textsc{Gauss} wrote to \textsc{Schumacher} on 30 May 1828: ‘which, I tell you, has probably anticipated \(\frac{1}{3}\) of my own investigations, and agrees with them in part even in the letters chosen to denote the quantities’.

The following passage from a letter of \textsc{Crelle} to \textsc{Abel} of 15 May 1828 probably refers to the same work: --- --- Here is what Mr.\ \textsc{Gauss} of Göttingen writes to me, whom I had also asked to send me something on elliptic functions, with which, as I have learned, he has been occupied for more than 30 years. ‘Other occupations prevent me for the moment from writing up these investigations. Mr.\ \textsc{Abel} has anticipated me by at least a third. He has just taken precisely the same route from which I set out in 1798. Thus I am not at all surprised that, for the greater part, he has arrived at the same results. Since, moreover, he has shown so much sagacity, penetration and elegance in his derivation, I consider myself thereby excused from writing up my own investigations.’ --- ---

\SourceLeaf{505}

We also possess a statement of \textsc{Gauss} about \textsc{Legendre}. In an undated letter he writes to \textsc{Olbers}: ‘In your last letter [probably the one marked ‘Bremen, 16 Aug.\ 1817, \textit{Received on 25 Aug.}’; the letters of those months are: O. to G., July 17. --- G. to O., Aug.\ 2. --- G. to O., the undated letter excerpted here --- O. to G., Nov.\ 2. --- G. to O., Dec.\ 2.], dearest friend, you asked for my opinion of \textsc{Mossotti}’s method of calculating the orbits of H. K., given in the Milan ephemerides. When I wrote to you recently, the subject was not sufficiently present to my mind, although I had earlier read that paper far enough to form a provisional opinion about it. At that moment my time did not allow me to study the matter properly again at once, and I therefore passed over your question. Since then I have had occasion to read that paper once more and enter further into its true spirit, and today I intend to spend an hour discussing this subject with you.’

‘Inclined, as I have always been, to welcome with affection every new original or inspired thought\HistoricalNote{31}, I was struck at my first reading by the truly new idea in \textsc{Mossotti}’s paper.’ --- ---

\medskip


\medskip
It particularly deserves to be stated that, concerning the theory of division of the lemniscatic arc, the manuscript estate contains nothing beyond the auxiliary propositions offered for it in the present printing, whereas in the work ‘\textit{Disquisitiones arithmeticae}’, issued in July 1801, Art.\ 335 of the seventh section, de aequationibus circuli sectiones definientibus, states: --- ‘Moreover, the principles of the theory which we are about to expound extend much more widely than they are carried here. For they can be applied not only to circular functions, but with equal success to many other transcendental functions, for example to those which depend on the integral \(\displaystyle\int\frac{\mathrm{d}x}{\surd(1-x^{4})}\), and furthermore also to various kinds of congruences. But since we are preparing a substantial separate work on those transcendental functions, and congruences will be treated at length in the continuation of the arithmetical investigations, it has seemed appropriate to consider only circular functions here.’ ---

\textsc{Abel} and \textsc{Jacobi} did not find \textsc{Gauss}’s investigations of elliptic functions available to them; they had to discover this field of science anew.

I shall attempt to present the particular relations between \textsc{Gauss}’s work in this field of pure analysis and the work of others in a separate publication, in connection with a history of \textsc{Gauss}’s entire scientific activity, whereas in these remarks appended to his own works I have included only the relevant documented facts.

\par Göttingen, June 1865.
\begin{flushright}
\textsc{Schering}.
\end{flushright}
\begin{center}
\rule{0.27\linewidth}{0.2pt}
\end{center}

\clearpage
\endgroup
